% experiments/exp2_chain_recipe/section_eval_callback.tex
%
% New short section that motivates a callback-aware standalone
% evaluation methodology for chain-recurrent memory and contrasts
% it with the standard chain eval (eval_chain.py) on the v14k
% step-3000 checkpoint of section7_v14_findings.tex.
%
% This file is meant to be \input'd after section8_recipe_revised.tex
% in the master draft. It does not modify draft.tex.

\section{Callback-aware standalone evaluation}
\label{sec:eval_cb}

The standalone evaluation pipeline used throughout
\S\S\ref{sec:experiments}--\ref{sec:revised_recipe}, as implemented
in \texttt{tools/eval\_chain.py}, scores cross-entropy uniformly
across a contiguous tail of the chain controlled by
\texttt{--eval\_window}. On D4v2 with the default
\texttt{eval\_window}${=}4$, the score window spans the callback
session plus the three preceding sessions, so the CE average
aggregates one callback session worth of \emph{meaningful} tokens
(the $\sim\!10$ \texttt{chain\_callback\_position} positions on
which the writer is supposed to help) against three sessions of
filler-and-evidence tokens whose ground-truth next-token
distribution is, by construction, invariant under whether memory
is recruited. The signal-to-noise of evidence-lift in this regime
is below one: a $+1.5$~nat per-token effect on $\sim\!10$ of
$\sim\!2000$ scored positions integrates to $\sim\!+0.0075$~nat
average, well below the bootstrap CI of $\dnm$ at few-hundred-chain
$n$. The standard chain eval is therefore not a measurement of
whether memory is helping the writer's stated job; it is a
measurement of per-token NLL averaged over a corpus that is mostly
not the writer's stated job.

\paragraph{The callback-aware standalone eval.} We define a
callback-aware variant that restricts the CE score to (a) the
callback session, identified by the chain's
\texttt{chain\_callback\_session\_idx} annotation, and (b) the
$\sim\!10$ token positions inside that session annotated by
\texttt{chain\_callback\_position}. Within this restricted token
set we report $\dnm^{\text{cb}}$, $\dsh^{\text{cb}}$, and
evidence-lift; each differs from \S\ref{sec:experiments} only in
the score-token mask. The $\dnm^{\text{cb}}$ floor is taken not
against $\Mc{=}\mathbf{0}$ --- which on a synthetic callback corpus
admits a trivial template solution --- but against
$\Mc^{\text{redact}}$, the chain's match memory recomputed with
the chain's evidence sessions replaced by distractor sessions
sampled from filler positions of other D4v2 chains. The redaction
preserves the recurrent update path while removing chain-specific
evidence, so $\Mc^{\text{redact}}$ exposes precisely the
chain-specific information the writer was supposed to compress.
The implementation is released at \texttt{tools/eval\_callback.py}.

\begin{table}[t]
\centering
\caption{Three eval methodologies applied to the v14k step-$3000$
checkpoint on the D4v2 validation split. The standard chain eval
averages one callback session of meaningful tokens against three
sessions of filler-and-evidence and integrates the signal to noise.
Restricting to the callback session alone does not recover the
effect: the callback session is itself $\sim\!90\%$ filler context
for the $\sim\!10$ \texttt{chain\_callback\_position} target tokens.
The callback-aware eval recovers a $+1.4$~nat memory-on-callback
effect. $n$ counts scored token-positions (top two rows) or chains
(bottom row).}
\label{tab:eval_methodology_compare}
\small
\begin{tabular}{lrrrr}
\toprule
methodology & $\dnm$ / $\dnm^{\text{cb}}$ & $\dsh$ / $\dsh^{\text{cb}}$ & evidence\_lift & $n$ \\
\midrule
standard (\texttt{score\_window}${=}4$)
                                          & $-0.085$ & $+0.001$ & $-0.085$ & $2000$ tok \\
callback-only (\texttt{score\_window}${=}1$)
                                          & $-0.095$ & $+0.003$ & $-0.095$ & $500$ tok \\
\textbf{callback-aware (this section)}    & $\boldsymbol{+1.44}$ & $\boldsymbol{+0.054}$ & $\boldsymbol{+0.071}$ & $200$ chains \\
\bottomrule
\end{tabular}
\end{table}

\paragraph{Discussion.} The callback-aware metric reveals a
$\sim\!1.5$~nat memory-on-callback effect on the v14k step-$3000$
checkpoint that the standard full-window metric averages to
\emph{negative} noise on the same checkpoint. The two metrics are
not in conflict: the standard metric averages CE over a window in
which $\sim\!99\%$ of scored tokens are not the writer's job, and
the recurrent readout adds a small per-token noise floor on those
tokens that integrates to $\dnm{=}{-0.085}$. The callback-aware
metric measures CE precisely on the tokens that are the writer's
job, against a baseline that admits no template shortcut, and
recovers the per-token effect at full magnitude. For
chain-recurrent memory architectures whose stated mechanism is
\emph{abstract callback}, callback-aware metrics should be the
default standalone-evaluation surface --- not because they are
systematically more favourable to memory architectures (a recipe
whose writer fails to recruit memory registers $\dnm^{\text{cb}}{=}0$
as cleanly as it registers $\dnm{=}0$), but because they are
diagnostic of where memory is supposed to help and silent on where
it is not. We adopt the callback-aware metrics as the primary
standalone-evaluation surface for the v15 series, with the standard
chain eval retained as a sanity check on aggregate per-token NLL.
